\documentclass{article}
\usepackage{mathtools} 
\usepackage{fontspec}
\usepackage[UTF8]{ctex}
\usepackage{amsthm}
\usepackage{mdframed}
\usepackage{xcolor}
\usepackage{amssymb}
\usepackage{amsmath}


% 定义新的带灰色背景的说明环境 zremark
\newmdtheoremenv[
  backgroundcolor=gray!10,
  % 边框与背景一致，边框线会消失
  linecolor=gray!10
]{zremark}{说明}

% 通用矩阵命令: \flexmatrix{矩阵名}{元素符号}{行数}{列数}
\newcommand{\flexmatrix}[4]{
  \[
  #1 = \begin{pmatrix}
    #2_{11}     & #2_{12}     & \cdots & #2_{1#4}   \\
    #2_{21}     & #2_{22}     & \cdots & #2_{2#4}   \\
    \vdots      & \vdots      & \ddots & \vdots     \\
    #2_{#31}    & #2_{#32}    & \cdots & #2_{#3#4}
  \end{pmatrix}
  \]
}

% 简化版命令（默认矩阵名为A，元素符号为a）: \quickmatrix{行数}{列数}
\newcommand{\quickmatrix}[2]{\flexmatrix{A}{a}{#1}{#2}}

\begin{document}
\title{注释 13.2}
\author{张志聪}
\maketitle

\begin{zremark}
  已知
  \begin{align*}
    \lim\limits_{n \to +\infty} a_n = a \\
    \lim\limits_{n \to +\infty} b_n = b
  \end{align*}
  则
  \begin{align*}
    \lim\limits_{n \to +\infty} |a_n - b_n| = |a - b|
  \end{align*}
\end{zremark}
\textbf{证明：}
\begin{itemize}
  \item 方法一

        设函数$f(x) = |x|$，该函数是初等函数，在$\mathbb{R}$上连续，
        所以在$a - b$处连续，且函数值为$f(a-b) = |a - b|$，又有
        \begin{align*}
          \lim\limits_{n \to +\infty} (a_n - b_n) = a - b
        \end{align*}
        由归结原理可知
        \begin{align*}
          \lim\limits_{n \to +\infty} |a_n - b_n| = f(a - b) = |a - b|
        \end{align*}

  \item 方法二

        提示：
        \begin{align*}
          ||a_n - b_n| - |a - b||
           & \leq |(a_n - b_n) - (a - b)| \\
           & \leq |a_n - a| + |b_n - b|
        \end{align*}
\end{itemize}

\begin{zremark}
  证明： 定理13.9 推论
\end{zremark}

\textbf{证明：}
对任意$x \in I$，由实数的稠密性可知，一定存在$a, b \in I$且$a < x < b$，
即$x \in [a, b]$，由于连续函数列$\{f_n(x)\}$在区间$I$上内闭一致收敛于$f(x)$，
由定义$\{f_n(x)\}$在$[a, b]$上一致收敛于$f(x)$。利用定理13.9可得
$f(x)$在$[a, b]$上连续，所以$f(x)$在$x$点处也连续。

\begin{zremark}
  证明：在定理13.11的条件下，还可推出在$[a, b]$上$f_n \rightrightarrows f \, (n \to \infty)$.
\end{zremark}
\textbf{证明：}
习题2

\begin{zremark}
  定理13.8在函数项级数上的推广。

  极限交换定理：$\sum u_n(x)$在$(a, x_0) \cup (x_0, b)$上一致收敛于$S(x)$，
  且对每个$n$，$\lim\limits_{x \to x_0} u_n(x) = a_n$，则
  $\lim\limits_{x \to x_0} \sum u_n(x)$和$\sum a_n$存在且相等，即
  \begin{align*}
    \lim\limits_{x \to x_0} \sum u_n(x) = \sum \lim\limits_{x \to x_0} u_n(x)
  \end{align*}
\end{zremark}
\textbf{证明：}
先证$\sum a_n$是收敛数项级数.
设$\sum u_n(x)$的部分和函数列为$\{S_n(x)\}$，由于
$\sum u_n(x)$一致收敛于$S_(x)$，
即$\{S_n(x)\}$一致收敛于$S(x)$，
由函数列一致收敛的柯西准则，故对任意$\epsilon > 0$，有$N$，
当$n \geq N$时，对任意正整数$p$和对一切$x \in (a, x_0) \cup (x_0, b)$，有
\begin{align*}
  |S_n(x) - S_{n + p}(x)| < \epsilon \\
  |u_{n + 1}(x) + \cdots + u_{n + p}(x)| < \epsilon
\end{align*}
从而
\begin{align*}
  |a_{n + 1} + \cdots + a_{n + p}| 
  < \lim\limits_{x \to x_0} |u_{n + 1}(x) + \cdots + u_{n + p}(x)| 
  \leq \epsilon
\end{align*}
由级数收敛的柯西准则可知$\sum a_n$是收敛级数。

设$\sum a_n = A$。再证明$\lim\limits_{x \to x_0} \sum u_n(x) = A$。

由于$\sum u_n(x)$一致收敛于$S(x)$以及$\sum a_n$收敛于A，因此对任意$\epsilon > 0$，
存在正数$N$，当$n > N$时，对任意$x \in (a, x_0) \cup (x_0, b)$，有
\begin{align*}
  |S_n(x) - S(x)| < \frac{\epsilon}{3}, \\
  |A_n - A| < \frac{\epsilon}{3}
\end{align*}
同时成立，其中$A_n = \sum \limits_{k = 1}^n a_k$。
特别地$n = N+1$，有
\begin{align*}
  |S_{N + 1}(x) - S(x)| < \frac{\epsilon}{3}, \\
  |A_{N + 1} - A| < \frac{\epsilon}{3}
\end{align*}
又$\lim\limits_{x \to x_0} S_{N + 1}(x) = A_{N + 1}$，
故存在$\delta > 0$，当$0 < |x - x_0| < \delta$时，
\begin{align*}
  |S_{N+1}(x) - A_{N + 1}| < \frac{\epsilon}{3}
\end{align*}
这样，当$x$满足$0 < |x - x_0| < \delta$时，
\begin{align*}
  |S(x) - A| & \leq |S(x) - S_{N + 1}(x)| + |S_{N + 1}(x) - A_{N + 1}| + |A_{N + 1} - A| \\
             & = \frac{\epsilon}{3} + \frac{\epsilon}{3} + \frac{\epsilon}{3} \\
             & = \epsilon
\end{align*}
即$\lim\limits_{x \to x_0} \sum u_n(x) = A$.

\begin{zremark}
  疑问：定理13.11 推论
\end{zremark}
\textbf{证明: }
个人感觉这个推论的前提条件少了一条，且在P38最下方也隐含的说了下：

$\{f_n\}$的每一项在$[a, b]$上有连续的导函数。

\end{document}